\documentclass[aps,preprint]{revtex4}%
\usepackage{amsfonts}%
\usepackage{amsmath}%
\setcounter{MaxMatrixCols}{30}%
\usepackage{amssymb}%
\usepackage{graphicx}
%TCIDATA{OutputFilter=latex2.dll}
%TCIDATA{Version=5.50.0.2953}
%TCIDATA{CSTFile=revtex4.cst}
%TCIDATA{Created=Monday, July 21, 2008 10:04:09}
%TCIDATA{LastRevised=Monday, July 21, 2008 10:39:29}
%TCIDATA{<META NAME="GraphicsSave" CONTENT="32">}
%TCIDATA{<META NAME="SaveForMode" CONTENT="1">}
%TCIDATA{BibliographyScheme=Manual}
%TCIDATA{<META NAME="DocumentShell" CONTENT="Articles\SW\REVTeX 4">}
%BeginMSIPreambleData
\providecommand{\U}[1]{\protect\rule{.1in}{.1in}}
%EndMSIPreambleData
\newtheorem{theorem}{Theorem}
\newtheorem{acknowledgement}[theorem]{Acknowledgement}
\newtheorem{algorithm}[theorem]{Algorithm}
\newtheorem{axiom}[theorem]{Axiom}
\newtheorem{claim}[theorem]{Claim}
\newtheorem{conclusion}[theorem]{Conclusion}
\newtheorem{condition}[theorem]{Condition}
\newtheorem{conjecture}[theorem]{Conjecture}
\newtheorem{corollary}[theorem]{Corollary}
\newtheorem{criterion}[theorem]{Criterion}
\newtheorem{definition}[theorem]{Definition}
\newtheorem{example}[theorem]{Example}
\newtheorem{exercise}[theorem]{Exercise}
\newtheorem{lemma}[theorem]{Lemma}
\newtheorem{notation}[theorem]{Notation}
\newtheorem{problem}[theorem]{Problem}
\newtheorem{proposition}[theorem]{Proposition}
\newtheorem{remark}[theorem]{Remark}
\newtheorem{solution}[theorem]{Solution}
\newtheorem{summary}[theorem]{Summary}
\newenvironment{proof}[1][Proof]{\noindent\textbf{#1.} }{\ \rule{0.5em}{0.5em}}
\begin{document}
\preprint{HEP/123-qed}
\title[Short title for running header]{Long title that appears on the title page}
\author{First Author}
\affiliation{My Institution}
\author{Second Author}
\affiliation{My Institution}
\author{Third Author}
\affiliation{Other Institution}
\keywords{one two three}
\pacs{PACS number}


%

\begin{align*}
\left\vert g_{+}\right\rangle  & =\frac{1}{\sqrt{2}}\left\{  \left\vert
\varphi_{1}\alpha\varphi_{3}\beta\right\rangle +\left\vert \varphi_{2}%
\alpha\varphi_{4}\beta\right\rangle \right\}  \\
& \equiv\frac{1}{\sqrt{2}}\left\{  \left\vert 1\alpha3\beta\right\rangle
+\left\vert 2\alpha4\beta\right\rangle \right\}
\end{align*}%
\[
\left\vert g_{-}\right\rangle =\frac{1}{\sqrt{2}}\left\{  \left\vert
1\alpha3\beta\right\rangle -\left\vert 2\alpha4\beta\right\rangle \right\}  ,
\]%
\begin{align*}
D^{1}  & =\frac{1}{4}\left\{  \left\vert 1\alpha\right\rangle \left\langle
1\alpha\right\vert +\left\vert 2\alpha\right\rangle \left\langle
2\alpha\right\vert \right.  \\
& \left.  +\left\vert 3\beta\right\rangle \left\langle 3\beta\right\vert
+\left\vert 4\beta\right\rangle \left\langle 4\beta\right\vert \right\}
\end{align*}%
\begin{align*}
D^{2}\left(  g_{+}\right)    & =\frac{1}{2}\left\{  \left\vert 1\alpha
3\beta\right\rangle \left\langle 1\alpha3\beta\right\vert +\left\vert
1\alpha3\beta\right\rangle \left\langle 2\alpha4\beta\right\vert \right.  \\
& \left.  +\left\vert 2\alpha4\beta\right\rangle \left\langle 1\alpha
3\beta\right\vert +\left\vert 2\alpha4\beta\right\rangle \left\langle
2\alpha4\beta\right\vert \right\}
\end{align*}%
\begin{align*}
D^{2}\left(  g_{-}\right)    & =\frac{1}{2}\left\{  \left\vert 1\alpha
3\beta\right\rangle \left\langle 1\alpha3\beta\right\vert -\left\vert
1\alpha3\beta\right\rangle \left\langle 2\alpha4\beta\right\vert \right.  \\
& \left.  -\left\vert 2\alpha4\beta\right\rangle \left\langle 1\alpha
3\beta\right\vert +\left\vert 2\alpha4\beta\right\rangle \left\langle
2\alpha4\beta\right\vert \right\}  .
\end{align*}%
\begin{align*}
&  \mathcal{E}\left(  g_{+}\right)  =\frac{1}{2}\left\{  \overset
{}{\left\langle 1\alpha\left\vert h_{1}1\alpha\right.  \right\rangle
+\left\langle 2\alpha\left\vert h_{1}2\alpha\right.  \right\rangle
+\left\langle 3\beta\left\vert h_{1}3\beta\right.  \right\rangle +\left\langle
4\beta\left\vert h_{1}4\beta\right.  \right\rangle }\right.  \\
&  \left.  \,+\left\langle 1\alpha3\beta\right\vert \left\vert 1\alpha
3\beta\right\rangle +\left\langle 2\alpha4\beta\right\vert \left\vert
2\alpha4\beta\right\rangle +\left\langle 1\alpha3\beta\right\vert \left\vert
2\alpha4\beta\right\rangle +\left\langle 2\alpha4\beta\right\vert \left\vert
1\alpha3\beta\right\rangle \overset{}{}\right\}
\end{align*}%
\begin{align*}
&  \mathcal{E}\left(  g_{-}\right)  =\frac{1}{2}\left\{  \overset
{}{\left\langle 1\alpha\left\vert h_{1}1\alpha\right.  \right\rangle
+\left\langle 2\alpha\left\vert h_{1}2\alpha\right.  \right\rangle
+\left\langle 3\beta\left\vert h_{1}3\beta\right.  \right\rangle +\left\langle
4\beta\left\vert h_{1}4\beta\right.  \right\rangle }\right.  \\
&  \left.  \,+\left\langle 1\alpha3\beta\right\vert \left\vert 1\alpha
3\beta\right\rangle +\left\langle 2\alpha4\beta\right\vert \left\vert
2\alpha4\beta\right\rangle -\left\langle 1\alpha3\beta\right\vert \left\vert
2\alpha4\beta\right\rangle -\left\langle 2\alpha4\beta\right\vert \left\vert
1\alpha3\beta\right\rangle \overset{}{}\right\}  .
\end{align*}%
\begin{align*}
\left\vert g_{\text{res}+}\right\rangle  &  =\left(  \sqrt{2}\right)
^{-1}\left\{  \left\vert 1\alpha1\beta\right\rangle +\left\vert 2\alpha
2\beta\right\rangle \right\}  \\
\left\vert g_{\text{res}-}\right\rangle  &  =\left(  \sqrt{2}\right)
^{-1}\left\{  \left\vert 1\alpha1\beta\right\rangle -\left\vert 2\alpha
2\beta\right\rangle \right\}  .
\end{align*}%
\[
\left\vert g\left(  \boldsymbol{c},\boldsymbol{\psi},\theta\right)
\right\rangle =c_{1}\left\vert 1\alpha3\beta\right\rangle +c_{2}e^{i\theta
}\left\vert 2\alpha4\beta\right\rangle ,
\]
where
\begin{align*}
&  c_{1}^{2}+c_{2}^{2}=1\\
&  0\leq\theta<2\pi.
\end{align*}
The constraint on the canonical coefficients leads to the parameterization
\begin{align*}
c_{1} &  =\cos\chi\\
c_{2} &  =\sin\chi\\
0 &  \leq\chi\leq\frac{\pi}{2}%
\end{align*}
and hence%
\[
\left\vert g\left(  \chi,\boldsymbol{\psi},\theta\right)  \right\rangle
=\cos\chi\left\vert 1\alpha3\beta\right\rangle +e^{i\theta}\sin\chi\left\vert
2\alpha4\beta\right\rangle .
\]
All the states with a fixed $\left\{  \chi,\boldsymbol{\psi}\right\}  $, where
$\boldsymbol{\psi\in}\left[  \boldsymbol{\psi}\right]  $, have a common FORDO
given by%
\[
D^{1}\left(  \chi,\boldsymbol{\psi}\right)  =\cos^{2}\chi\left\{  \left\vert
1\alpha\right\rangle \left\langle 1\alpha\right\vert +\left\vert
3\beta\right\rangle \left\langle 3\beta\right\vert \right\}  +\sin^{2}%
\chi\left\{  \left\vert 2\alpha\right\rangle \left\langle 2\alpha\right\vert
+\left\vert 4\beta\right\rangle \left\langle 4\beta\right\vert \right\}
\]
but different SORDOs%
\begin{align*}
D^{2}\left(  \chi,\boldsymbol{\psi}\right)   &  =\cos^{2}\chi\left\vert
1\alpha3\beta\right\rangle \left\langle 1\alpha3\beta\right\vert +\sin^{2}%
\chi\left\vert 2\alpha4\beta\right\rangle \left\langle 2\alpha4\beta
\right\vert +\\
&  \frac{1}{2}\left\{  e^{-i\theta}\left\vert 1\alpha3\beta\right\rangle
\left\langle 2\alpha4\beta\right\vert +e^{i\theta}\left\vert 2\alpha
4\beta\right\rangle \left\langle 1\alpha3\beta\right\vert \right\}  \sin2\chi.
\end{align*}
These states have different energies that can be expressed by%
\begin{align*}
&  \mathcal{E}\left(  \theta\right)  =\left\{  \left\langle 1\alpha\left\vert
h_{1}1\alpha\right.  \right\rangle +\left\langle 3\beta\left\vert h_{1}%
3\beta\right.  \right\rangle \right\}  \cos^{2}\chi+\\
&  \left\{  \left\langle 2\alpha\left\vert h_{1}2\alpha\right.  \right\rangle
+\left\langle 4\beta\left\vert h_{1}4\beta\right.  \right\rangle \right\}
\sin^{2}\chi\\
&  +\left\langle 1\alpha3\beta\right\vert \left\vert 1\alpha3\beta
\right\rangle \cos^{2}\chi+\left\langle 2\alpha4\beta\right\vert \left\vert
2\alpha4\beta\right\rangle \sin^{2}\chi\\
&  +\frac{1}{2}\left\{  e^{i\theta}\left\langle 1\alpha3\beta\right\vert
\left\vert 2\alpha4\beta\right\rangle +e^{-i\theta}\left\langle 2\alpha
4\beta\right\vert \left\vert 1\alpha3\beta\right\rangle \right\}  \sin2\chi.
\end{align*}
After spin integration eq. $\left(  \ref{EnTheta}\right)  $ gives
\begin{align*}
&  \mathcal{E}\left(  \theta\right)  =\left\{  \left(  1\left\vert
h_{1}1\right.  \right)  +\left(  3\left\vert h_{1}3\right.  \right)  \right\}
\cos^{2}\chi+\left\{  \left(  2\left\vert h_{1}2\right.  \right)  +\left(
4\left\vert h_{1}4\right.  \right)  \right\}  \sin^{2}\chi\\
&  \,+\left(  \left.  13\right\vert 13\right)  \cos^{2}\chi+\left(  \left.
24\right\vert 24\right)  \sin^{2}\chi+\operatorname{Re}\left\{  e^{i\theta
}\left(  \left.  13\right\vert 24\right)  \sin2\chi\right\}  \\
&  \equiv k\left(  \chi,\boldsymbol{\psi}\right)  +f\left(  \chi
,\boldsymbol{\psi},\theta\right)  ,
\end{align*}
where
\[
\left(  \left.  jk\right\vert lm\right)  =\int\varphi_{j}^{\ast}\left(
\boldsymbol{r}_{1}\right)  \varphi_{k}^{\ast}\left(  \boldsymbol{r}%
_{2}\right)  \frac{1}{\left\Vert \boldsymbol{r}_{1}-\boldsymbol{r}%
_{2}\right\Vert }\varphi_{l}\left(  \boldsymbol{r}_{1}\right)  \varphi
_{m}\left(  \boldsymbol{r}_{2}\right)  d\boldsymbol{r}_{1}d\boldsymbol{r}_{2},
\]%
\begin{align*}
&  \mathcal{E}\left(  \theta\right)  =\left\{  \left(  1\left\vert
h_{1}1\right.  \right)  +\left(  3\left\vert h_{1}3\right.  \right)  \right\}
\cos^{2}\chi+\left\{  \left(  2\left\vert h_{1}2\right.  \right)  +\left(
4\left\vert h_{1}4\right.  \right)  \right\}  \sin^{2}\chi\\
&  \,+\left(  \left.  13\right\vert 13\right)  \cos^{2}\chi+\left(  \left.
24\right\vert 24\right)  \sin^{2}\chi+\operatorname{Re}\left\{  e^{i\theta
}\left(  \left.  13\right\vert 24\right)  \sin2\chi\right\}  \\
&  \equiv k\left(  \chi,\boldsymbol{\psi}\right)  +f\left(  \chi
,\boldsymbol{\psi},\theta\right)  ,
\end{align*}
where
\[
\left(  \left.  jk\right\vert lm\right)  =\int\varphi_{j}^{\ast}\left(
\boldsymbol{r}_{1}\right)  \varphi_{k}^{\ast}\left(  \boldsymbol{r}%
_{2}\right)  \frac{1}{\left\Vert \boldsymbol{r}_{1}-\boldsymbol{r}%
_{2}\right\Vert }\varphi_{l}\left(  \boldsymbol{r}_{1}\right)  \varphi
_{m}\left(  \boldsymbol{r}_{2}\right)  d\boldsymbol{r}_{1}d\boldsymbol{r}_{2},
\]
%

\[
\left(  j\left\vert h_{1}j\right.  \right)  =\int\varphi_{j}^{\ast}\left(
\boldsymbol{r}\right)  \left\{  -\frac{1}{2}\nabla_{\boldsymbol{r}}%
^{2}+V_{\text{ext}}\left(  \boldsymbol{r}\right)  \right\}  \varphi_{j}\left(
\boldsymbol{r}\right)  d\boldsymbol{r},
\]%
\begin{align*}
&  k\left(  \chi,\boldsymbol{\psi}\right)  =\left\{  \left(  1\left\vert
h_{1}1\right.  \right)  +\left(  3\left\vert h_{1}3\right.  \right)  \right\}
\cos^{2}\chi\\
&  +\left\{  \left(  2\left\vert h_{1}2\right.  \right)  +\left(  4\left\vert
h_{1}4\right.  \right)  \right\}  \sin^{2}\chi\\
&  +\left(  \left.  13\right\vert 13\right)  \cos^{2}\chi+\left(  \left.
24\right\vert 24\right)  \sin^{2}\chi
\end{align*}
and
\[
f\left(  \chi,\boldsymbol{\psi},\theta\right)  =\operatorname{Re}\left\{
e^{i\theta}\left(  \left.  13\right\vert 24\right)  \sin2\chi\right\}  .
\]
Over the domain of such geminals one can define a DMF as
\[
\mathcal{F}\left(  \chi,\boldsymbol{\psi}\right)  =\mathcal{F}\left(
D^{1}\left(  \chi,\boldsymbol{\psi}\right)  \right)  =\min_{\boldsymbol{\theta
}}\mathcal{E}\left(  \theta\right)  .
\]
The only part of the functional $\mathcal{E}$ that depends on $\theta$ is
$f\left(  \chi,\boldsymbol{\psi},\theta\right)  ,$ which one can express as
\begin{align*}
f\left(  \chi,\boldsymbol{\psi},\theta\right)    & =\operatorname{Re}\left\{
e^{i\theta}\left(  \left.  13\right\vert 24\right)  \right\}
=\operatorname{Re}\left\{  e^{i\theta}ae^{i\alpha}\right\}  \\
& =\operatorname{Re}\left\{  e^{i\left(  \theta+\alpha\right)  }a\right\}
=a\cos\left(  \theta+\alpha\right)  ,
\end{align*}
where
\begin{align*}
a &  =\left\vert \left(  \left.  13\right\vert 24\right)  \right\vert \\
\alpha &  =\arg\left\{  \left(  \left.  13\right\vert 24\right)  \right\}  .
\end{align*}
The minimum of this function can be easily evaluated giving
\begin{align*}
f\left(  \chi,\boldsymbol{\psi},\theta_{\min}\right)   &  =\min_{\theta
}\left\{  a\cos\left(  \theta+\alpha\right)  \right\}  =-a\\
\theta_{\min} &  =\pi-\alpha.
\end{align*}
Hence the value of the DMF, $\mathcal{F},$ at the point $D^{1}\left(
\chi,\boldsymbol{\psi}\right)  $ is%
\[
\mathcal{F}\left(  \chi,\boldsymbol{\psi}\right)  =k\left(  \chi
,\boldsymbol{\psi}\right)  -2\left\vert \left(  \left.  13\right\vert
24\right)  \right\vert \sin2\chi
\]
and the minimum of the AGP energy functional $\mathcal{E}\left(
\boldsymbol{c}\left(  \chi\right)  ,\boldsymbol{\psi},\theta\right)  ,$ which
determines the AGP\ ground state, is given by
\[
\mathcal{E}\left(  \left\vert g_{GS}\right\rangle \right)  =\min
_{\chi,\boldsymbol{\psi}}\mathcal{F}\left(  \chi,\boldsymbol{\psi}\right)  .
\]



\end{document}
